% ============================================================
% Section 5: Conclusions and Future Research
% ============================================================

\section{Conclusions and Future Research}
\label{sec:conclusions}

This work presents an implementation of Physics-Informed Neural Networks with
sinusoidal activation (PINN-SIREN) for solving the 2D Helmholtz equation with a
complex field, under a \textit{mesh-free} paradigm
\cite{Karniadakis2021_review, Cuomo2022_review}.
The two completed experiments (NB01 and NB02) validate the architecture with high
accuracy.

The main conclusions are:

\begin{enumerate}
  \item \textbf{Exceptional accuracy in 1D and 2D.}
    The SIREN architecture reached $L^2$ errors of $0.006\%$ (1D) and $0.171\%$ (2D),
    both several orders of magnitude below the threshold adopted in this work
    ($< 5\%$).

  \item \textbf{The Adam + L-BFGS strategy is essential.}
    Adam provides robust global exploration; L-BFGS refines the solution using
    second-order curvature information.
    Using either optimizer alone produces inferior results.

  \item \textbf{The physics weight $\lambda = 0.1$ balances data and physics in 2D.}
    In the 2D case, the physical loss function includes two residuals (real and
    imaginary), effectively doubling the weight of the physics if $\lambda = 1.0$ is
    used.
    Calibrating to $\lambda = 0.1$ was critical for convergence.
\end{enumerate}

The natural extensions of this work include:

\begin{itemize}
  \item \textbf{Optical speckle simulation} (NB03): applying the model validated in
    NB02 to speckle generation, replacing the plane-wave condition with a
    random-phase boundary $\phi(x) \sim \mathcal{U}(0, 2\pi)$ at $y = 0$, and
    statistically validating the generated pattern with Goodman's contrast
    criterion $|C - 1| < 0.1$.

  \item \textbf{FEM benchmark} (NB05): quantitative comparison against FEniCSx for
    multiple speckle realizations, quantifying the \textit{Speed-up Factor}
    $S = T_{\mathrm{FEM}} / T_{\mathrm{PINN}}$ \cite{Thuerey2021_PBDL} and the peak
    memory footprint, in line with the objectives of this research.

  \item \textbf{Inhomogeneous media} (future work, notebook numbering not yet
    confirmed): extension to the Helmholtz equation with a variable refractive index
    $n(\mathbf{r})$, relevant for speckle in GRIN materials and biological tissues.
    Postponed as agreed with the thesis advisor.
\end{itemize}

\section*{Acknowledgments}

Null

\section*{Usage of Generative AI}

% NOTE FOR THE AUTHOR: adjust this statement so that it accurately reflects the
% actual use of generative AI before submitting the manuscript.
NUll
